% Перечень условных обозначений оформляется в алфавитном порядке
\sectioncentered*{Перечень условных обозначений}
\addcontentsline{toc}{section}{Перечень условных обозначений}
\label{sec:reduction}

API — Application Programming Interface, программный интерфейс приложения;

CNN — Convolutional Neural Network, свёрточная нейронная сеть;

CV — Computer Vision, компьютерное зрение;

DGT — Digital Game Technology, производитель электронного оборудования для регистрации шахматных партий;

FEN — Forsyth--Edwards Notation, стандарт текстового представления позиции на шахматной доске;

FPS — Frames Per Second, количество кадров в секунду;

FSM — Finite State Machine, конечный автомат;

HUD — Head-Up Display, графический информационный слой поверх видеопотока;

JSONL — JSON Lines, построчный формат хранения структурированных записей JSON;

MJPEG — Motion JPEG, формат потоковой передачи видеопоследовательности;

ML — Machine Learning, машинное обучение;

NNUE — Efficiently Updatable Neural Network, нейросетевая оценочная функция шахматного движка;

ONNX — Open Neural Network Exchange, открытый формат представления моделей машинного обучения;

PGN — Portable Game Notation, стандарт текстового представления шахматной партии;

RANSAC — Random Sample Consensus, алгоритм робастной оценки параметров модели при наличии выбросов;

REST — Representational State Transfer, архитектурный стиль взаимодействия по протоколу HTTP;

RTSP — Real Time Streaming Protocol, потоковый протокол передачи видеоданных;

SAN — Standard Algebraic Notation, стандартная алгебраическая шахматная нотация;

SHA-256 — Secure Hash Algorithm 256-bit, криптографическая хеш-функция с длиной дайджеста 256 бит;

SSIM — Structural Similarity Index Measure, индекс структурного сходства изображений;

UCI — Universal Chess Interface, протокол взаимодействия шахматного программного обеспечения с движком;

ViT — Vision Transformer, архитектура трансформера для обработки изображений;

WebSocket — протокол полнодуплексной связи поверх TCP для обновления состояния интерфейса в реальном времени;

YOLO — You Only Look Once, семейство методов детектирования объектов.
